\subsection{Factorization through the constitutive operator}

Retain the lifted angular chart of
Section~\ref{iii:sec:hojas}. To make the angular conversion explicit,
write
\begin{equation}
 A_{\deg}=\frac{180x}{\pi},\qquad
 C^*_{\deg}=\frac{180y}{\pi},\qquad
 q_\pm=e^{-x\mp y},\qquad s_\pm=q_\pm^{30}.
 \label{iii:eq:barbero-angular-lift}
\end{equation}
The orientation considered in \eqref{iii:eq:rchannels} corresponds to
$0<y<x$; sheet interchange is expressed by $y\mapsto-y$ in the
larger domain $x>|y|$. Real exponentials act on these
lifts, which preserve the values of $x\pm y$ during
composition.

The oriented Barbero--Immirzi functional is
\begin{equation}
 \gamma=\frac{\pi A_{\deg}C^*_{\deg}}{180}
   +12\bigl[S_{90}(q_-)-S_{90}(q_+)\bigr]
   -\bigl[S_{120}(q_-)-S_{120}(q_+)\bigr].
 \label{iii:eq:barbero-functional}
\end{equation}
The first component is the angular evaluation; the second is the
difference of the fundamental-chain readout on the conjugate channels.
The arguments are received from the previously constructed angular transport.
The formula determines a dimensionless scalar on that starting
structure.

Let $\psi:(3/4,1)\to(0,1)$ be the inverse of $f$ constructed through
\eqref{iii:eq:cubic}. Define
\begin{equation}
 \mathcal H(s)=\frac{12s^3}{1-s^9}
             -\frac{s^4}{1-s^{12}}
             -\frac{(\log s)^2}{20\pi},\qquad 0<s<1.
 \label{iii:eq:barbero-h}
\end{equation}
This function is real analytic on its domain. The exponents $3,4,9,12$
result from expressing heights $90,120$ and returns $270,360$ in
units of $30$; the logarithmic coefficient is fixed when
the angular contribution is transported.

\begin{theorem}[Oriented constitutive expression of Barbero--Immirzi]
\label{iii:thm:barbero-factorization}
In chart \eqref{iii:eq:barbero-angular-lift},
\begin{equation}
 \gamma=\mathcal H(s_-)-\mathcal H(s_+)
       =\mathcal H\bigl(\psi(r_-)\bigr)
        -\mathcal H\bigl(\psi(r_+)\bigr).
 \label{iii:eq:barbero-factorization}
\end{equation}
On the two-dimensional channel representation,
\begin{equation}
 \gamma=-\operatorname{Tr}
    \left[\mathcal R\,\mathcal H\bigl(\psi(\mathcal V)\bigr)\right],
 \label{iii:eq:barbero-trace}
\end{equation}
where $\operatorname{Tr}$ is the ordinary trace, with
$\operatorname{Tr}I=2$.
\end{theorem}
\begin{proof}
Substitution of $s=q^{30}$ into \eqref{iii:eq:barbero-return-series} gives
\[
 S_{90}(q)=\frac{s^3}{1-s^9},\qquad
 S_{120}(q)=\frac{s^4}{1-s^{12}}.
\]
Moreover, \eqref{iii:eq:barbero-angular-lift} implies
\[
 \log s_\pm=-30(x\pm y)
 =-\frac\pi6(A_{\deg}\pm C^*_{\deg}).
\]
By the difference of squares,
\begin{align}
 \frac{(\log s_+)^2-(\log s_-)^2}{20\pi}
 &=\frac{900\bigl[(x+y)^2-(x-y)^2\bigr]}{20\pi}\\
 &=\frac{180xy}{\pi}
 =\frac{\pi A_{\deg}C^*_{\deg}}{180}.
 \label{iii:eq:barbero-log-term}
\end{align}
Adding these identities proves the first equality in
\eqref{iii:eq:barbero-factorization}. The second uses the inversion
$\psi(f(s_\pm))=s_\pm$ proved in
Section~\ref{iii:sec:constitutiva}.

For the operator expression, the sheet projectors have rank one and
$\mathcal RP_\pm=\pm P_\pm$. Spectral calculus gives
\[
 \mathcal H\bigl(\psi(\mathcal V)\bigr)
 =\mathcal H(s_+)P_++\mathcal H(s_-)P_-.
\]
Multiplying by $\mathcal R$ and taking the ordinary trace gives
$\mathcal H(s_+)-\mathcal H(s_-)$. Its negative is $\gamma$.
\end{proof}

The orientation mark $\mathcal R$ belongs to the functional. If that
mark is held fixed and the channels of $\mathcal V$ are interchanged,
the sign of $\gamma$ changes. Evaluation of the spectrum without its labels preserves
the unordered pair, whereas \eqref{iii:eq:barbero-trace} preserves the
oriented contribution. If the normalized trace
$\tau_2=\operatorname{Tr}/2$ is used, the same formula reads
$\gamma=-2\tau_2[\mathcal R\mathcal H(\psi(\mathcal V))]$.

By contrast, a simultaneous change of representation
$(\mathcal R,\mathcal V)\mapsto
(U\mathcal RU^{-1},U\mathcal VU^{-1})$ preserves the functional:
spectral calculus is transported by conjugation and the trace is
invariant. Interchanging the response relative to a fixed orientation
and jointly transporting response and orientation are therefore
different operations.

This normalization is compatible with the amplifications of
\eqref{iii:eq:logvac}. For
$\mathcal V_m=\mathcal V\otimes I_{9^m}$ and
$\mathcal R_m=\mathcal R\otimes I_{9^m}$,
\begin{equation}
 \gamma=-\frac1{9^m}\operatorname{Tr}
 \left[\mathcal R_m\mathcal H\bigl(\psi(\mathcal V_m)\bigr)\right].
 \label{iii:eq:barbero-amplification}
\end{equation}
Indeed, functional calculus preserves the tensor product and the trace
of $I_{9^m}$ is $9^m$. Representation multiplicity is separated
from the oriented evaluation.

\subsection{Reconstruction from vacuum coordinates}

Identities \eqref{iii:eq:four-identities} express the same
functional through normalized impedance and velocity:
\begin{align}
 \gamma={}&
 \mathcal H\!\left(\psi\!\left(
     \sqrt{\frac{\widehat Z}{\widehat c}}\right)\right)
 -\mathcal H\!\left(\psi\!\left(
     \frac1{\sqrt{\widehat Z\widehat c}}\right)\right).
 \label{iii:eq:barbero-zc}
\end{align}
The positive domain of this reconstruction is described by
\begin{equation}
 1<\widehat Z\widehat c<\frac{16}{9},\qquad
 \frac9{16}<\frac{\widehat Z}{\widehat c}<1.
 \label{iii:eq:barbero-zc-domain}
\end{equation}
Indeed, these inequalities are equivalent to
$3/4<r_\pm<1$ after substitution of the positive roots of
\eqref{iii:eq:four-identities}. For orientation $y>0$, one also retains
$r_-<r_+$, equivalently $\widehat Z<1$. Domain points
are used with the angular compatibility already established by
the TPK; inversion is a subsequent reconstruction of its channels.
The coordinates are the internal ones of \eqref{iii:eq:four}: a change of
units must be undone through \eqref{iii:eq:undo-units} before
$\psi$ is applied.

Interchange $y\mapsto-y$ acts as
$(\widehat Z,\widehat c)\mapsto(\widehat Z^{-1},\widehat c)$ and
produces $\gamma\mapsto-\gamma$. In the balanced case $y=0$,
the channels coincide and \eqref{iii:eq:barbero-factorization} gives
$\gamma=0$. Normalized velocity preserves the product of the
channels; its joint readout with impedance recovers both
responses and the orientation required to evaluate the functional.
In this balanced case $\mathcal V=rI$ has degenerate spectrum:
the mark $\mathcal R$ is retained as representation data, although
the scalar $r$ no longer distinguishes its projectors. For simple spectrum,
by contrast,
$P_+=(\mathcal V-r_-I)/(r_+-r_-)$ and $P_-=I-P_+$.
